\documentclass[11pt]{article}

\usepackage[T1]{fontenc}
\usepackage{amsmath,amssymb,amsthm}
\usepackage[margin=1in]{geometry}
\usepackage[hidelinks]{hyperref}
\newcommand{\sdp}{\mathbin{\times\!\!|}}
\newcommand{\wprod}{\mathbin{\operatorname{wr}}}
\hypersetup{
  pdftitle={The Disconnected Graph A(4,4,3): Classification and a Counterexample to Pan's arXiv-v1 Conjecture},
  pdfauthor={Xiaojian Zeng}
}

\newtheorem{theorem}{Theorem}
\DeclareMathOperator{\Aut}{Aut}
\newcommand{\Cay}{\operatorname{Cay}}
\newcommand{\Sn}{S_n}
\newcommand{\Sfour}{S_4}
\newcommand{\Afour}{A_4}
\newcommand{\Vfour}{V_4}

\title{The Disconnected Graph $A(4,4,3)$:\\
Classification and a Counterexample to Pan's arXiv-v1 Conjecture}
\author{Xiaojian Zeng\\
PBC School of Finance, Tsinghua University, 100190\\
\texttt{xiaojianzeng@foxmail.com}}
\date{}

\begin{document}

\maketitle

\begin{abstract}
For $F_t$ the set of permutations of $S_n$ having exactly $t$ fixed
points, Pan conjectured in arXiv:2407.19745v1 that the full automorphism
group of $\Cay(S_n,F_t)$ is
$[R(S_n)\sdp\operatorname{Inn}(S_n)]\sdp C_2$ for every admissible
$n$ and $t$.  We determine the graph and the full group at
$(n,t)=(4,1)$.  The connection set is $F_1=A_4\setminus V_4$, and its
cosets give
\[
  \Cay(S_4,F_1)\cong A(4,4,3)\cong 2K_{4,4,4}.
\]
Consequently
\[
  \Aut(A(4,4,3))\cong (S_4\wprod S_3)\wprod S_2,
\]
of order $13{,}759{,}414{,}272$, whereas the conjectured group has order
$1{,}152$.  Thus the conjecture as stated in that arXiv version is false
at this parameter.  No assertion is made about priority or the text of
any later journal version.
\end{abstract}

\section{Context and conventions}

Day and Tripathi introduced arrangement graphs as a flexible class of
interconnection networks \cite{DayTripathi1992}.  Chen, Ghorbani, and Wong
then considered the generalized graph $A(n,k,r)$: its vertices are the
ordered $k$-tuples of distinct elements of $[n]$, and two vertices are
adjacent when they differ in exactly $r$ coordinates
\cite{ChenGhorbaniWong2014}.  Pan determined the full automorphism groups
in two extreme families and proposed a uniform formula for the associated
symmetric-group Cayley graphs \cite{Pan2024}.  The corresponding full
groups for the complete derangement and complete transposition graphs
were determined by Deng and Zhang \cite{DengZhang2011} and Ganesan
\cite{Ganesan2015}, respectively.

We address only Conjecture~3.1 of Pan's arXiv version~1 and only the
parameter $(n,t)=(4,1)$.  We use Pan's left-Cayley convention: vertices
$g,h\in S_4$ are adjacent when
\[
 g\sim h \quad\Longleftrightarrow\quad hg^{-1}\in F_1.
 \tag{1}\label{eq:cayley}
\]
Permutations act as functions on $[4]$, and $g$ is identified with the
tuple $[g(1),g(2),g(3),g(4)]$.  The symbol $\sdp$ denotes a semidirect
product; wreath products below use their natural
imprimitive actions, and $2K_{4,4,4}$ denotes the disjoint union of two
copies of $K_{4,4,4}$.

\section{Classification and full automorphism group}

Set
\[
 \Vfour=\{1,(12)(34),(13)(24),(14)(23)\}.
\]

\begin{theorem}\label{thm:main}
Let $F_1$ be the permutations of $S_4$ with exactly one fixed point.  Then
\[
 F_1=A_4\setminus V_4,
 \qquad
 \Cay(S_4,F_1)\cong A(4,4,3)\cong 2K_{4,4,4},
\]
and
\[
 \Aut(A(4,4,3))\cong (S_4\wprod S_3)\wprod S_2.
\]
The group has order $13{,}759{,}414{,}272$.  In particular,
Conjecture~3.1 of \cite{Pan2024}, as stated in arXiv version~1, is false
at $(n,t)=(4,1)$.
\end{theorem}

\begin{proof}
A permutation of four letters with exactly one fixed point acts without a
fixed point on the other three letters, so it is a $3$-cycle.  Conversely,
every $3$-cycle fixes the remaining letter.  Hence
\[
 F_1=\{(123),(132),(124),(142),(134),(143),(234),(243)\}.
\]
The even permutations in $S_4$ consist of the identity, the three double
transpositions, and these eight $3$-cycles.  Therefore
\[
 F_1=A_4\setminus V_4. \tag{2}\label{eq:connection}
\]
The displayed $V_4$ is normal in $S_4$: conjugation preserves the cycle
type of its three nonidentity elements.

We first check the arrangement-graph identification without changing the
Cayley convention.  For $g,h\in S_4$, a coordinate $i$ satisfies
$h(i)=g(i)$ exactly when $g(i)$ is a fixed point of $hg^{-1}$.  The map
$i\mapsto g(i)$ makes this a bijection between agreeing coordinates and
fixed points.  Thus \eqref{eq:cayley} holds exactly when the two
permutation tuples differ in three coordinates.  This proves
$\Cay(S_4,F_1)\cong A(4,4,3)$.

By \eqref{eq:cayley} and \eqref{eq:connection},
\begin{align*}
 g\sim h
 &\Longleftrightarrow hg^{-1}\in A_4\setminus V_4\\
 &\Longleftrightarrow
   \operatorname{sgn}(g)=\operatorname{sgn}(h)
   \quad\text{and}\quad V_4g\ne V_4h. \tag{3}\label{eq:cosets}
\end{align*}
Indeed, when $g$ and $h$ have the same parity, $hg^{-1}\in A_4$, while
$hg^{-1}\in V_4$ is equivalent to $h\in V_4g$, and hence to equality of
the two $V_4$-cosets.  Normality means these right cosets are also left
cosets.

Each $V_4$-coset has four elements of one parity.  Each of the twelve
even and twelve odd permutations therefore splits into three such
cosets.  Equation~\eqref{eq:cosets} says that within a parity class there
are no edges inside a coset and all possible edges between distinct
cosets.  Each parity class consequently induces $K_{4,4,4}$.  It is
connected, and there are no edges between the two parity classes, so
these classes are precisely the two connected components.  Hence
\[
 A(4,4,3)\cong 2K_{4,4,4}. \tag{4}\label{eq:decomp}
\]

Let $Y=K_{4,4,4}$, with parts $P_1,P_2,P_3$.  Every independent set of
$Y$ lies in one part, so its inclusion-maximal independent sets are
exactly $P_1,P_2,P_3$.  Every automorphism therefore permutes the three
parts.  The kernel of this action consists of the arbitrary, independent
permutations inside them, and every permutation of the parts preserves
adjacency.  The action splits, giving
\[
 \Aut(Y)\cong (S_4)^3\sdp S_3=S_4\wprod S_3,
 \qquad
 |\Aut(Y)|=(4!)^3 3!=82{,}944. \tag{5}\label{eq:onecomponent}
\]
By \eqref{eq:decomp}, an automorphism of the full graph permutes its two
connected components.  The kernel is $\Aut(Y)^2$, and an exchange of the
two copies supplies the nontrivial element of $S_2$.  Therefore
\[
 \Aut(A(4,4,3))
 \cong \Aut(Y)^2\sdp S_2
 \cong (S_4\wprod S_3)\wprod S_2, \tag{6}
\]
whose order is
\[
 (82{,}944)^2\,2=13{,}759{,}414{,}272. \tag{7}
\]

Finally, $S_4$ has trivial center, so
$|\operatorname{Inn}(S_4)|=24$.  The iterated semidirect product proposed
in \cite[Conjecture~3.1]{Pan2024} has order
\[
 |R(S_4)|\,|\operatorname{Inn}(S_4)|\,|C_2|
 =24\cdot24\cdot2=1{,}152. \tag{8}
\]
The orders in (7) and (8) differ.  Since the cited conjecture quantifies
over every admissible $n$ and $t$, the case $(4,1)$ disproves it.
\end{proof}

\section{Exact check, scope, and assistance disclosure}

As a corroborating check, a Python standard-library program enumerated
all $24$ permutation tuples and reconstructed adjacency directly from
the connection set.  Using only tuple and integer operations, it verified
the eight elements of $F_1$, degree $8$, two parity components of order
$12$, three four-element $V_4$-cosets in each component, all within- and
between-part adjacencies, and the three group-order formulas displayed in
\eqref{eq:onecomponent}, (7), and (8).  The proof of
Theorem~\ref{thm:main} does not depend on this enumeration.  This
separation between a mathematical
argument and a reconstructive exact checker is also used in finite graph
certificate work where the verifier rebuilds the graph and its symmetry
action before accepting stored data \cite[Sections~2 and~6]{ZengEtAl2026}.

The conclusion is deliberately restricted to the conjecture in
arXiv:2407.19745v1.  It does not claim priority and makes no statement
about whether a later journal text retains, modifies, or omits the
conjecture.

OpenAI Codex assisted with literature searching and source verification,
the preparation of the exact checker, adversarial organization of the
argument, and initial manuscript drafting.  Neither language-model output
nor the finite check is used as a premise in the proof above.  No proof
assistant was used.

\raggedright
\bibliographystyle{plain}
\bibliography{references}

\end{document}
